\section{相关工作}

\subsection{语言模型推理与自改进}

Chain-of-thought（CoT）提示证明了显式中间推理步骤对复杂问题求解的重要性~\citep{wei2022chain_cot}。在此基础上，Synthetic Prompting、STaR、Tree of Thoughts、Quiet-STaR 等方法进一步探索了如何自动生成、搜索和利用推理轨迹~\citep{shao2023synthetic_synprompt, zelikman2022star, yao2023tree_tot, zelikman2024quiet_quietstar}。与此同时，自一致性、自验证、自反思和基于检索的纠错机制也被广泛用于提升推理可靠性~\citep{wang2022self_consistency, weng2023large_self_verification, madaan2023self_selfrefine, li2025cot_rag, niu2024mitigating, zheng2023ddcot}。

不过，这些方法大多面向单轮推理或可直接验证的任务；对于 GUI 这种长轨迹、多步执行、环境状态持续变化的场景，单纯依赖模型自我解释并不足够可靠。CSRS 的贡献就在于把步骤级 reasoning 放到轨迹级成功信号之下进行校准，从而让“自改进”真正适用于 GUI 智能体。

\subsection{GUI 智能体}

近期大量工作开始系统研究 GUI 智能体的数据、训练和架构设计，包括 UI-TARS、OpenCUA、GUI-Owl、UITron、Mobile-Agent-v3 等~\citep{qin2025ui_uitars1, wang2025ui_uitars2, wang2025opencua, ye2025mobile_mobileagentv3, zeng2025uitron}。与此同时，Claude、Seed、Qwen 等通用基础模型也开始原生支持 GUI 场景~\citep{claudeopus45, guo2025seed1, Bai2025Qwen3VL}。

大多数 GUI 智能体遵循 ReAct 范式：观察、推理、再行动~\citep{yao2022react}。最近的方法则进一步引入多智能体协作、分层规划以及视觉与代码混合控制~\citep{liu2025pc, song2025coact1, zhang2025ufo2, agashe2025agents2}。Step-GUI 的不同之处在于：它不仅关注模型架构和执行流程，还围绕“如何稳定、低成本地构造高质量训练数据”建立了完整闭环。

\subsection{GUI 评测基准}

GUI 评测大致可以分成两类。一类是静态基准，例如 ScreenSpot、ScreenSpot-v2、ScreenSpot-Pro、FuncPred、MoTIF、RefExp、LlamaTouch 和 VisualWebBench 等，主要考察元素定位和动作预测~\citep{cheng2024seeclick, wu2024atlas, li2025screenspotpro, li2025autogui, burns2022dataset_motif, rastogi2021uibert_refexp, zhang2024llamatouch, liu2024visualwebbench_vwb}。另一类是交互式基准，例如 AndroidWorld 与 OSWorld，它们允许智能体在更真实的系统环境中执行完整任务~\citep{rawles2024androidworld, xie2024osworld}。

AndroidDaily 与这些工作互补：它更强调真实日常高频移动使用场景，特别是中文移动生态中的出行、购物与本地生活任务，因此能更直接反映 GUI 智能体走向真实部署时的能力边界。
